\documentclass[12pt]{article}
\usepackage[utf8]{inputenc}
\usepackage[T1]{fontenc}
\usepackage[margin=1in]{geometry}
\usepackage{setspace}
\usepackage{xcolor}
\usepackage[colorlinks=true, urlcolor=blue!60!black, linkcolor=blue!60!black]{hyperref}
\setlength{\parindent}{0em}
\setlength{\parskip}{0.6em}
\renewcommand{\baselinestretch}{1.0}
\begin{document}
\begin{center}
{\LARGE \textit{Rust v. Sullivan}, 500 U.S. 173 (1991)\par}
\vspace{0.5em}
{\large Case Brief --- Two-Page Edition\par}
\end{center}

\section{Memory Jogger}
The Title X ``gag rule'' case: government may define the scope of a program it funds, barring funded projects from abortion counseling, referral, or advocacy (\href{https://supreme.justia.com/cases/federal/us/500/173/}{500 U.S. 173}).

\section{Detailed Case Facts}
Title X funds family planning projects but bars funds from ``programs where abortion is a method of family planning.'' After 18 years of permitting abortion counseling, HHS in 1988 forbade Title X projects to counsel or refer for abortion, to advocate it, or to share facilities with abortion activities. Grantees, physicians including Dr. Irving Rust, and New York sued.

\section{Procedural History}
The district court and a divided Second Circuit upheld the regulations; the First and Tenth Circuits disagreed. The Supreme Court affirmed.

\section{Judicial Votes}
Affirmed, 5--4: Rehnquist wrote for White, Kennedy, Scalia, and Souter; Blackmun, Stevens, and O'Connor each dissented, with Marshall joining Blackmun.

\section{Holding}
The regulations are a permissible reading of an ambiguous statute under \textit{Chevron}. The government may fund one activity to the exclusion of another and limit speech within the funded project, so long as grantees remain free to speak outside it; refusing to subsidize abortion counseling does not burden a woman's due process rights.

\section{Analysis of Opinions}
Relying on \textit{Maher}, \textit{Harris v. McRae}, and \textit{Regan}, Rehnquist held that selective funding is not viewpoint discrimination and distinguished conditions on a \textit{grantee}, which may not reach its private speech, from conditions on a \textit{project}. He cautioned that public forums and universities may be different. Blackmun called the rules viewpoint-based intrusions into the physician--patient relationship; Stevens read \S~1008 to reach only conduct; O'Connor would have avoided the constitutional question on statutory grounds.

\section{Examples: Future Applications}
\textbf{Same-Side:} a smoking-cessation grant may bar funded counselors from recommending vaping. \textbf{Opposite-Side:} requiring grantees to adopt an organization-wide anti-prostitution pledge reaches speech outside the program (\textit{AID v. AOSI}). \textbf{Fence-Sitter:} speech conditions on university research grants, given academic freedom.

\section{Critique}
Defenders stress democratic control over public money; libertarians warn that conditional grants let government buy silence. Progressive scholars such as Robert Post and David Cole criticize the co-opting of doctors into delivering state-scripted messages and the unstable line between \textit{Rust} and \textit{Velazquez}.

\section{Key Quotations}
\begin{quote}
``The Government can, without violating the Constitution, selectively fund a program to encourage certain activities it believes to be in the public interest, without at the same time funding an alternate program which seeks to deal with the problem in another way. In so doing, the Government has not discriminated on the basis of viewpoint; it has merely chosen to fund one activity to the exclusion of the other.''
\end{quote}
\noindent --- Chief Justice Rehnquist, Opinion of the Court, 500 U.S. at 193.

\end{document}
